\documentclass[12pt]{article}
\usepackage[a4paper,margin=1in]{geometry}
\usepackage{amsmath,amssymb,bm}

\title{Double-Link Pendulum Dynamics (Tasks 2.1 and 2.2)}
\author{}
\date{}

\begin{document}
\maketitle

\section*{Given data}
\[
m_1 = 1\,\text{kg},\quad m_2 = 1\,\text{kg},\quad
l_1 = 0.8\,\text{m},\quad l_2 = 1.2\,\text{m},\quad
g = 9.8\,\text{m/s}^2.
\]

We choose generalized coordinates
\[
\bm q = \begin{bmatrix}\alpha_1 \\ \alpha_2\end{bmatrix},
\]
where $\alpha_1$ and $\alpha_2$ are the link angles measured from the positive $x$ axis.
Links are massless rods and point masses are concentrated at their ends. Motor inertias are neglected.

\section*{2.1 Euler-Lagrange derivation}
The point-mass positions are
\[
\bm r_1 =
\begin{bmatrix}
l_1\cos\alpha_1 \\
l_1\sin\alpha_1
\end{bmatrix},
\qquad
\bm r_2 =
\begin{bmatrix}
l_1\cos\alpha_1 + l_2\cos\alpha_2 \\
l_1\sin\alpha_1 + l_2\sin\alpha_2
\end{bmatrix}.
\]

Velocities:
\[
\dot{\bm r}_1 =
\begin{bmatrix}
-l_1\dot\alpha_1\sin\alpha_1 \\
l_1\dot\alpha_1\cos\alpha_1
\end{bmatrix},
\quad
\dot{\bm r}_2 =
\begin{bmatrix}
-l_1\dot\alpha_1\sin\alpha_1 - l_2\dot\alpha_2\sin\alpha_2 \\
l_1\dot\alpha_1\cos\alpha_1 + l_2\dot\alpha_2\cos\alpha_2
\end{bmatrix}.
\]
Hence,
\[
\|\dot{\bm r}_1\|^2 = l_1^2\dot\alpha_1^2,
\]
\[
\|\dot{\bm r}_2\|^2 = l_1^2\dot\alpha_1^2 + l_2^2\dot\alpha_2^2
+ 2l_1l_2\dot\alpha_1\dot\alpha_2\cos(\alpha_1-\alpha_2).
\]

Kinetic energy:
\[
T = \frac12 m_1\|\dot{\bm r}_1\|^2 + \frac12 m_2\|\dot{\bm r}_2\|^2
= \frac12 (m_1+m_2)l_1^2\dot\alpha_1^2
+ \frac12 m_2l_2^2\dot\alpha_2^2
+ m_2l_1l_2\dot\alpha_1\dot\alpha_2\cos(\alpha_1-\alpha_2).
\]

Potential energy (with $y$ positive upward):
\[
V = m_1gy_1 + m_2gy_2
= (m_1+m_2)gl_1\sin\alpha_1 + m_2gl_2\sin\alpha_2.
\]

Lagrangian:
\[
\mathcal L = T - V.
\]

Using
\[
\frac{d}{dt}\left(\frac{\partial\mathcal L}{\partial \dot\alpha_i}\right)
- \frac{\partial\mathcal L}{\partial \alpha_i} = 0,
\qquad i=1,2,
\]
the equations of motion are
\begin{align*}
&(m_1+m_2)l_1^2\ddot\alpha_1
+ m_2l_1l_2\cos(\alpha_1-\alpha_2)\ddot\alpha_2
+ m_2l_1l_2\sin(\alpha_1-\alpha_2)\dot\alpha_2^2
+ (m_1+m_2)gl_1\cos\alpha_1 = 0,\\
&m_2l_1l_2\cos(\alpha_1-\alpha_2)\ddot\alpha_1
+ m_2l_2^2\ddot\alpha_2
- m_2l_1l_2\sin(\alpha_1-\alpha_2)\dot\alpha_1^2
+ m_2gl_2\cos\alpha_2 = 0.
\end{align*}

Substituting the given numerical values:
\begin{align*}
&1.28\,\ddot\alpha_1 + 0.96\cos(\alpha_1-\alpha_2)\,\ddot\alpha_2
+ 0.96\sin(\alpha_1-\alpha_2)\,\dot\alpha_2^2
+ 15.68\cos\alpha_1 = 0,\\
&0.96\cos(\alpha_1-\alpha_2)\,\ddot\alpha_1 + 1.44\,\ddot\alpha_2
- 0.96\sin(\alpha_1-\alpha_2)\,\dot\alpha_1^2
+ 11.76\cos\alpha_2 = 0.
\end{align*}

\section*{2.2 Matrix form $\bm M(\bm q)\ddot{\bm q}+\bm C(\bm q,\dot{\bm q})\dot{\bm q}+\bm g(\bm q)=\bm 0$}
Let $\Delta = \alpha_1-\alpha_2$. Then
\[
\bm M(\bm q)=
\begin{bmatrix}
(m_1+m_2)l_1^2 & m_2l_1l_2\cos\Delta \\
m_2l_1l_2\cos\Delta & m_2l_2^2
\end{bmatrix},
\]
\[
\bm C(\bm q,\dot{\bm q})=
\begin{bmatrix}
0 & m_2l_1l_2\sin\Delta\,\dot\alpha_2 \\
-m_2l_1l_2\sin\Delta\,\dot\alpha_1 & 0
\end{bmatrix},
\]
\[
\bm g(\bm q)=
\begin{bmatrix}
(m_1+m_2)gl_1\cos\alpha_1 \\
m_2gl_2\cos\alpha_2
\end{bmatrix}.
\]

Numerically,
\[
\bm M(\bm q)=
\begin{bmatrix}
1.28 & 0.96\cos(\alpha_1-\alpha_2) \\
0.96\cos(\alpha_1-\alpha_2) & 1.44
\end{bmatrix},
\]
\[
\bm C(\bm q,\dot{\bm q})=
\begin{bmatrix}
0 & 0.96\sin(\alpha_1-\alpha_2)\,\dot\alpha_2 \\
-0.96\sin(\alpha_1-\alpha_2)\,\dot\alpha_1 & 0
\end{bmatrix},
\]
\[
\bm g(\bm q)=
\begin{bmatrix}
15.68\cos\alpha_1 \\
11.76\cos\alpha_2
\end{bmatrix}.
\]

Therefore,
\[
\bm M(\bm q)\ddot{\bm q}+\bm C(\bm q,\dot{\bm q})\dot{\bm q}+\bm g(\bm q)=\bm 0.
\]

\end{document}